%document
\documentclass[10pt]{beamer}
%theme
\usetheme{metropolis}
% packages
\usepackage{color}
\usepackage{listings}
\usepackage[ngerman]{babel}
\usepackage[utf8x]{inputenc}
\usepackage{multicol}


% color definitions
\definecolor{mygreen}{rgb}{0,0.6,0}
\definecolor{mygray}{rgb}{0.5,0.5,0.5}
\definecolor{mymauve}{rgb}{0.58,0,0.82}

\lstset{
    backgroundcolor=\color{white},
    % choose the background color;
    % you must add \usepackage{color} or \usepackage{xcolor}
    basicstyle=\footnotesize\ttfamily,
    % the size of the fonts that are used for the code
    breakatwhitespace=false,
    % sets if automatic breaks should only happen at whitespace
    breaklines=true,                 % sets automatic line breaking
    captionpos=b,                    % sets the caption-position to bottom
    commentstyle=\color{mygreen},    % comment style
    % deletekeywords={...},
    % if you want to delete keywords from the given language
    extendedchars=true,
    % lets you use non-ASCII characters;
    % for 8-bits encodings only, does not work with UTF-8
    frame=single,                    % adds a frame around the code
    keepspaces=true,
    % keeps spaces in text,
    % useful for keeping indentation of code
    % (possibly needs columns=flexible)
    keywordstyle=\color{blue},       % keyword style
    % morekeywords={*,...},
    % if you want to add more keywords to the set
    numbers=left,
    % where to put the line-numbers; possible values are (none, left, right)
    numbersep=5pt,
    % how far the line-numbers are from the code
    numberstyle=\tiny\color{mygray},
    % the style that is used for the line-numbers
    rulecolor=\color{black},
    % if not set, the frame-color may be changed on line-breaks
    % within not-black text (e.g. comments (green here))
    stepnumber=1,
    % the step between two line-numbers.
    % If it's 1, each line will be numbered
    stringstyle=\color{mymauve},     % string literal style
    tabsize=4,                       % sets default tabsize to 4 spaces
    % show the filename of files included with \lstinputlisting;
    % also try caption instead of title
    language = Java,
	showspaces = false,
	showtabs = false,
	showstringspaces = false,
	escapechar = ,
}

\def\ContinueLineNumber{\lstset{firstnumber=last}}
\def\StartLineAt#1{\lstset{firstnumber=#1}}
\let\numberLineAt\StartLineAt



\newcommand{\codeline}[1]{
	\alert{\texttt{#1}}
}

% This Document contains the information about this course.

% Authors of the slides
\author{Fritz Windisch, Lukas Tassanyi}

% Name of the Course
\institute{Java-Course}

% Fancy Logo
\titlegraphic{\hfill\includegraphics[height=1.25cm]{../templates/fsr_logo_cropped}}


\usepackage{qtree}

\title{Java}
\subtitle{Exceptions}
\date{\today}

\begin{document}

\begin{frame}
\titlepage
\end{frame}

\begin{frame}{Overview}
\tableofcontents
\end{frame}

\section{Exceptions}
\subsection{Overview}
\begin{frame}{}
    While running software many things can go wrong. 
    You have to deal with errors or exceptional behavior. %AE
    \vfill
    Java offers exception handling out of the box.
    Exceptions seperate error-handling from normal code.
    \vfill
    On this slide \emph{exception} means the Java term and \emph{error} a nonspecified general term.
\end{frame}

\begin{frame}{Hierarchy}
    \Tree [.Object [.Throwable Error [.Exception \dots{} RuntimeException ] ] ] \\
    \vfill
    Every exception is a subclass of \emph{Throwable}. 
    \emph{Error} is also a subclass of \emph{Throwable} but used for serious errors
    like \emph{VirtualMachineError}. \\
    
    \scriptsize\url{https://docs.oracle.com/javase/7/docs/api/java/lang/Throwable.html}
\end{frame}

\begin{frame}{Checked Exceptions}
    Every exception except \emph{RuntimeException} and its subclasses are \textbf{checked exceptions}.
    \vfill
    A checked exception has to be handled or denoted.
    \vfill
    The cause of this kind of exception is often outside of your program.
\end{frame}

\begin{frame}{Unchecked Exceptions}
    \emph{RuntimeException} and its subclasses are called \textbf{unchecked exceptions}.
    \vfill
    Unchecked Exceptions do not have to be denoted or handled, but can be.
    Often handling is senseless because the program can not recover 
    in case such exception occurs.
    \vfill
    The cause of an unchecked exception can be a method call with incorrect arguments.
    Therefore any method could throw an unchecked exception.
    Most unchecked exceptions are caused by the programer.
    \vfill
    Errors are also unchecked.
\end{frame}

\subsection{Catching Exceptions}    
\begin{frame}[fragile]{Introduction}
    \begin{lstlisting}[basicstyle=\ttfamily\scriptsize]
    public class Calc {
    
        public static void main(String[] args) {
    
            int a = 7 / 0;
            // will cause an ArithmeticException
            
            System.out.println(a);
        }
    }
    \end{lstlisting}
    A division by zero causes an \emph{ArithmeticException} which is a subclass of \emph{RuntimeException}. 
    Therefore \emph{ArithmeticException} is unchecked and does not have to be handled.
\end{frame}

\begin{frame}[fragile]{Try and Catch}
    Nevertheless the exception can be handled.
    \begin{lstlisting}[basicstyle=\ttfamily\scriptsize]
    public class Calc {
    
        public static void main(String[] args) {
    
            try {
                int a = 7 / 0;
            } catch (ArithmeticException e) {
                System.out.println("Division by zero.");
            }
        }
    }
    \end{lstlisting}
    The \textbf{catch}-block, also called exception handler,
    is invoked if the specified exception (ArithmeticException) occurs in the \textbf{try}-block.\\
    In general there can be multiple catch-blocks handling multiple kinds of exceptions.
\end{frame}

\begin{frame}[fragile]{Stack Trace}
    \begin{lstlisting}[basicstyle=\ttfamily\scriptsize]
    public class Calc {
    
        public static void main(String[] args) {
    
            try {
                int a = 7 / 0;
            } catch (ArithmeticException e) {
                System.out.println("Division by zero.");
                e.printStackTrace();
            }
        }
    }
    \end{lstlisting}
    The stack trace shows the order of method calls leading to point where the exception occurs.
\end{frame}

\begin{frame}[fragile]{Stack Trace}
    \begin{lstlisting}[basicstyle=\ttfamily\scriptsize]
    Division by zero.
    java.lang.ArithmeticException: / by zero
        at Calc.main(Calc.java:6)
        at sun.reflect.NativeMethodAccessorImpl.invoke0(Native Method)
        at sun.reflect.NativeMethodAccessorImpl.invoke(NativeMethodAccessorImpl.java:62)
        at sun.reflect.DelegatingMethodAccessorImpl.invoke(DelegatingMethodAccessorImpl.java:43)
        at java.lang.reflect.Method.invoke(Method.java:498)
        at com.intellij.rt.execution.application.AppMain.main(AppMain.java:147)
    \end{lstlisting}
\end{frame}

\begin{frame}[fragile]{Finally}
    \begin{lstlisting}[basicstyle=\ttfamily\scriptsize]
    public class Calc {
    
        public static void main(String[] args) {
    
            try {
                int a = 7 / 0;
            } catch (ArithmeticException e) {
                System.out.println("Division by zero.");
                e.printStackTrace();
            } finally {
                System.out.println("End of program.");
            }
        }
    }
    \end{lstlisting}
    The \textbf{finally}-block will always be executed, regardless if an exception occurs.
\end{frame}

\subsection{Throwing Exceptions}
\begin{frame}[fragile]{Propagate Exceptions}
    Unhandled exceptions can be thrown (propagated).
    \begin{lstlisting}[basicstyle=\ttfamily\scriptsize]    
    public static int divide (int divident, int divisor) throws ArithmeticException {
        return divident / divisor;
    }
    \end{lstlisting}
    The method \texttt{int divide(\dots)} propagates the exception to the calling
    method denoted by the keyword \textbf{throws}.
\end{frame}

\begin{frame}[fragile]{Propagate Exceptions - Test 1}
    \begin{lstlisting}[basicstyle=\ttfamily\scriptsize]
    public class Calc {
    
        public static int divide (int divident, int divisor) throws ArithmeticException {
            return divident / divisor;
        }
    
        public static void main(String[] args) {
    
            int a = 0;
            try {
                a = Calc.divide(7, 0);
            } catch (ArithmeticException e) {
                System.out.println("Division by zero.");
                e.printStackTrace();
            }
        }
    }
    \end{lstlisting}
\end{frame}

\begin{frame}[fragile]{Propagate Exceptions - Test 2}
    \begin{lstlisting}[basicstyle=\ttfamily\scriptsize, firstnumber=7]
    public static void main(String[] args) {
    
        int a = 0;
        try {
            a = Calc.divide(7, 0);
        } catch (ArithmeticException e) {
            System.out.println("Division by zero.");
            e.printStackTrace();
        }
    }
    \end{lstlisting}
    In this example there are two jumps in the stack trace:\\
    \texttt{java.lang.ArithmeticException: / by zero}\\
    \texttt{at Calc.divide(Calc.java:4)}\\
    \texttt{at Calc.main(Calc.java:11)}
\end{frame}

\begin{frame}{Java API}
    The Java API shows\footnote{\scriptsize\url{https://docs.oracle.com/javase/7/docs/api/}}
    if a method throws exceptions. 
    The notation \texttt{throws exception} means that the method can throw 
    exceptions in case of an unexpected situation.
    It does not mean that the method throws exception every time.
    \vfill
    Check if the Exception is a subclass of \emph{RuntimeException}. 
    If not the exception has to be handled or rethrown.
\end{frame}

\begin{frame}[fragile]{Creating new Exceptions}
    You can create und use your own exception class.
    \begin{lstlisting}[basicstyle=\ttfamily\scriptsize]
    public class DivisionByZeroException extends Exception {

    }
    \end{lstlisting}
    \vfill
    \begin{lstlisting}[basicstyle=\ttfamily\scriptsize]
    public static int divide (int divident, int divisor) throws DivisionByZeroException {
        if (divisor == 0) {
            throw new DivisionByZeroException();
        }
        return divident / divisor;
    }
    \end{lstlisting}    
    Exceptions can be thrown manually with the keyword \textbf{throw}.
\end{frame}

\begin{frame}[fragile]{Creating new Exceptions - Test}
    \begin{lstlisting}[basicstyle=\ttfamily\scriptsize]
    public static void main(String[] args) {
    
        int a = 0;
        try {
            a = Calc.divide(7, 0);
        } catch (DivisionByZeroException e) {
            System.out.println("Division by zero.");
            e.printStackTrace();
        }
    }
    \end{lstlisting}    
    \emph{DivisionByZeroException} is checked and therefore has to be handled.
\end{frame}

\end{document}